% ============================================================
%  LATEX PATCH v2 — "Fungal Hyphae Reorganize Condensation Fields
%                     as Distributed Hygroscopic Sinks"
%  Generated: 2026-04-19
%
%  Updates for new figure structure:
%    Figure 3: panels A--E only (universal collapse)
%    Figure 4: NEW morphological analysis figure (panels A--F)
%
%  Numbers from: Final Results (FFT, local density CV, lacunarity)
% ============================================================


% ============================================================
% BLOCK 1 — FIGURE 3 CAPTION (panels A--E only)
%
% Figure 3 now shows universal collapse of fungi + hydrogels.
% All morphology panels (old G--J) removed; now in Figure 4.
% ============================================================

\begin{figure}[H]
\centering
\includegraphics[width=\textwidth]{Figure3.pdf}
\caption{\textbf{Universal condensation response across hydrogels and fungi.}
\textbf{(A)}~Representative brightfield images of three fungal species
on the cooled substrate: \textit{Aspergillus} (left),
\textit{Mucor} (center), and \textit{Rhizopus} (right).
Scale bar: 2~mm.
\textbf{(B)}~Normalized radial size profiles $R^* = R/R_\infty$
versus distance from the source boundary for all seven conditions
(4~hydrogel concentrations + 3~fungal species; 5~replicates each,
$n = 35$~trials). Shaded bands: mean~$\pm$~SEM.
\textbf{(C)}~Heatmap of $R^*$ for all 35~trials, sorted by
condition. Red bar: dry-zone width~$\delta$.
\textbf{(D)}~Half-attenuation distance $d^*$ versus~$\delta$
for all conditions ($R^2 = 0.90$).
Fungal species follow the same linear relationship established
by the hydrogel calibration series.
\textbf{(E)}~Normalized size gradient
$(R_\mathrm{far} - R_\mathrm{near})/R_\mathrm{mid}$
versus~$\delta$ ($R^2 = 0.92$).}
\label{fig:collapse}
\end{figure}


% ============================================================
% BLOCK 2 — FIGURE 4 CAPTION (NEW morphological analysis figure)
%
% Panels:
%   A: Aspergillus representative images (3 rows: colony surface,
%      light microscopy, Hessian tubeness overlay)
%   B: Mucor representative images (same 3 rows)
%   C: FFT spectral slope boxplot (colony surface, n=14+7)
%   D: FFT spectral slope boxplot (3D macro ROI, n=13+10)
%   E: Local density CV boxplot (light microscopy, n=6+3)
%   F: Lacunarity boxplot (3D macro, n=13+11)
% ============================================================

\begin{figure}[H]
\centering
\includegraphics[width=\textwidth]{Figure4.pdf}
\caption{\textbf{Multi-scale structural morphology distinguishes
\textit{Aspergillus} from \textit{Mucor}.}
\textbf{(A,\,B)}~Representative images of \textit{Aspergillus}~(A)
and \textit{Mucor}~(B) at three imaging scales.
Top row: colony surface micrographs showing natural packing
(scale bars: 1~mm).
Middle row: light micrographs of disaggregated tissue showing
unit morphology---dense conidial aggregates
(\textit{Aspergillus}) versus a distributed filamentous network
(\textit{Mucor}; scale bars: 1~mm).
Bottom row: multi-scale Hessian tubeness response (inferno
colormap), highlighting filamentary structures detected by
second-derivative eigenvalue analysis.
\textbf{(C)}~Spectral slope~$\alpha$ from tile-based 2D~FFT
of colony-surface micrographs.
\textit{Aspergillus} ($\alpha = -2.96 \pm 0.02$, $n = 14$)
exhibits a shallower slope than \textit{Mucor}
($\alpha = -3.09 \pm 0.10$, $n = 7$;
$p = 0.013$, Cohen's $d = 2.21$), indicating greater
fine-scale surface complexity.
\textbf{(D)}~Spectral slope~$\alpha$ from 2D~FFT of
manually selected in-focus regions on 3D colony-surface
photographs, replicating the spectral separation on an
independent image set
(\textit{Aspergillus}: $-3.22 \pm 0.09$, $n = 13$;
\textit{Mucor}: $-3.46 \pm 0.24$, $n = 10$;
$p = 0.012$, $d = 1.41$).
\textbf{(E)}~Local density coefficient of variation (CV) from
light microscopy of disaggregated tissue.
Foreground fraction was measured in 256~$\times$~256~\textmu m
patches; the CV across patches quantifies spatial
heterogeneity of tissue distribution.
\textit{Aspergillus} ($0.69 \pm 0.11$) is significantly more
heterogeneous than \textit{Mucor} ($0.39 \pm 0.03$;
$p = 7.0 \times 10^{-4}$, $d = 3.12$;
Benjamini--Hochberg adjusted $p = 0.012$).
The ratio is preserved at $10\times$, $20\times$, and
$40\times$ magnification ($1.7\times$--$1.9\times$),
confirming that heterogeneity reflects intrinsic morphology.
\textbf{(F)}~Lacunarity at the largest measurement scale,
computed by the gliding-box algorithm on Otsu-segmented 3D
colony-surface ROIs.
\textit{Mucor} ($2.04 \pm 0.54$, $n = 11$) has significantly
higher lacunarity than \textit{Aspergillus}
($1.31 \pm 0.16$, $n = 13$;
$p = 0.001$, $d = -1.92$), indicating a more
heterogeneous gap distribution at the macro scale.}
\label{fig:morphology}
\end{figure}


% ============================================================
% BLOCK 3 — RESULTS: Structural morphology paragraph
%
% Insert after the sentence ending "...separates from
%   \textit{Mucor} ($140 \pm 33$~\textmu m) and \textit{Rhizopus}
%   ($110 \pm 23$~\textmu m)."
%
% Replaces everything through "(INSERT SEM SECTION)".
% All references now point to Fig 4 (not Fig 3 G--J).
% ============================================================

This correspondence suggests that colony-scale architecture, rather
than cell-wall chemistry alone, modulates how far the depletion zone
extends.  To test this hypothesis we quantified the structural
morphology of \textit{Aspergillus} and \textit{Mucor} colonies at
three imaging scales using four independent metrics
(Fig.~\ref{fig:morphology}).

Automated two-dimensional Fourier analysis of colony-surface
micrographs revealed that \textit{Aspergillus} surfaces have a
shallower spectral slope ($\alpha = -2.96 \pm 0.02$, $n = 14$~images)
than \textit{Mucor} ($\alpha = -3.09 \pm 0.10$, $n = 7$; Welch's
$t = 3.4$, $p = 0.013$; Cohen's $d = 2.21$;
Fig.~\ref{fig:morphology}C), indicating greater power at high spatial
frequencies and therefore more fine-scale surface texture.  This
separation was replicated on an independent set of 3D colony-surface
photographs with manually selected in-focus regions
($\alpha_{\mathrm{Asp}} = -3.22 \pm 0.09$, $n = 13$;
$\alpha_{\mathrm{Muc}} = -3.46 \pm 0.24$, $n = 10$; $p = 0.012$,
$d = 1.41$; Fig.~\ref{fig:morphology}D).

Light microscopy of disaggregated colony tissue confirmed that
the two genera differ in unit morphology
(Fig.~\ref{fig:morphology}A,\,B).  \textit{Aspergillus} tissue
consists of dense conidial aggregates separated by voids, whereas
\textit{Mucor} tissue forms a distributed filamentous network.
The coefficient of variation of local patch density---a
preparation-independent metric of spatial heterogeneity---was
$0.69 \pm 0.11$ for \textit{Aspergillus} versus
$0.39 \pm 0.03$ for \textit{Mucor}
($p = 7.0 \times 10^{-4}$, $d = 3.12$;
Benjamini--Hochberg adjusted $p = 0.012$;
Fig.~\ref{fig:morphology}E), with the ratio preserved across
$10\times$, $20\times$, and $40\times$ magnification
($1.7\times$--$1.9\times$).

Finally, lacunarity analysis of Otsu-segmented 3D colony-surface
ROIs showed that \textit{Mucor} has a significantly more
heterogeneous gap distribution at the macro scale
($\Lambda = 2.04 \pm 0.54$, $n = 11$) than \textit{Aspergillus}
($\Lambda = 1.31 \pm 0.16$, $n = 13$; $p = 0.001$,
$d = -1.92$; Fig.~\ref{fig:morphology}F).  The combination of
high micro-scale heterogeneity (elevated CV) with low macro-scale
lacunarity is the signature of hierarchical porosity:
\textit{Aspergillus} is internally structured at fine scales but
fills the projected area uniformly, whereas \textit{Mucor}'s
tissue is locally uniform but irregularly distributed.


% ============================================================
% BLOCK 4 — METHODS: Spectral analysis
% (replaces old spectral analysis paragraph)
% ============================================================

\paragraph{Spectral analysis of colony surface morphology.}
Colony-surface micrographs were acquired using a DSLR macro
imaging system at matched magnification (${\sim}1$~\textmu m/px,
calibrated from a stage micrometer with 200~\textmu m divisions).
Two datasets were analyzed: (i)~14~\textit{Aspergillus} and
7~\textit{Mucor} colony-surface images (``colony surface''), and
(ii)~13~\textit{Aspergillus} and 10~\textit{Mucor} manually
selected in-focus regions from 3D colony photographs
(``3D macro''; pixel calibration 0.94~\textmu m/px from a
200~\textmu m-division grid).

Grayscale images were divided into overlapping tiles
($512 \times 512$~px, stride 256~px for colony surface;
$256 \times 256$~px, stride 128~px for 3D macro ROIs).
Tile-level quality control rejected tiles below the 10th
percentile of within-image Laplacian variance (out-of-focus) and
tiles with intensity standard deviation below 5~grey levels
(textureless); typical retention was 75--90\%.  For each retained
tile, a 2D Hann window was applied, the discrete 2D Fourier
transform was computed, and the magnitude-squared power spectrum
was azimuthally averaged into radial frequency bins.  A mean
radial power profile was assembled per image by averaging across
retained tiles.  The log--log spectral slope~$\alpha$ was
obtained by ordinary least-squares regression of
$\log_{10}(\mathrm{power})$ versus
$\log_{10}(\mathrm{frequency})$ over 0.01--0.45~cyc/px.
Because $\alpha$ is the exponent of a power-law fit in frequency
space, it is independent of absolute pixel calibration.
Statistical comparison used Welch's $t$-test.


% ============================================================
% BLOCK 5 — METHODS: Light microscopy (local density CV)
% ============================================================

\paragraph{Light microscopy of colony structure.}
Colony tissue was excised from mature cultures, disaggregated on
glass slides, and imaged under brightfield illumination at
$10\times$, $20\times$, and $40\times$ magnification.  Pixel size
was calibrated from a stage micrometer with 10~\textmu m divisions
(0.75~\textmu m\,px$^{-1}$ at $10\times$;
0.38~\textmu m\,px$^{-1}$ at $20\times$;
0.19~\textmu m\,px$^{-1}$ at $40\times$).  Grayscale images were
segmented by Otsu thresholding of the inverted intensity, followed
by morphological opening and closing.  Spatial heterogeneity was
quantified as the coefficient of variation of foreground-area
fraction measured across a regular grid of non-overlapping
$256 \times 256$~\textmu m patches.  Six~\textit{Aspergillus}
(two per magnification) and three~\textit{Mucor} (one per
magnification) fields were analyzed.  Benjamini--Hochberg
correction was applied to control the false-discovery rate across
all tested metrics.


% ============================================================
% BLOCK 6 — METHODS: Surface fragmentation and lacunarity
% ============================================================

\paragraph{Surface fragmentation analysis.}
Three-dimensional colony-surface photographs were cropped to
manually selected in-focus regions of interest
(13~\textit{Aspergillus}, 11~\textit{Mucor} ROIs; pixel
calibration 0.94~\textmu m/px).  Each ROI was binarized by Otsu
thresholding (dark~=~tissue) and cleaned by morphological opening
($3 \times 3$~kernel, 2~iterations) and closing
($5 \times 5$~kernel, 2~iterations).  Lacunarity was computed
using the gliding-box algorithm: for each box size~$r$
(4, 8, 16, 32, 64, 128, 256~px), the variance-to-mean ratio
of the foreground mass within sliding windows was recorded;
the reported value is the lacunarity at the largest box size
that fit within the ROI.  Statistical comparison used Welch's
$t$-test and Mann--Whitney $U$~test.


% ============================================================
% BLOCK 7 — DISCUSSION: Hierarchical porosity and Kelvin argument
%
% Insert after the existing discussion of δ differences.
% This replaces the old one-sentence spectral interpretation.
% ============================================================

\paragraph{Reproductive morphology as a determinant of sink strength.}
The stronger depletion-zone response of \textit{Aspergillus}
relative to \textit{Mucor} (Fig.~\ref{fig:collapse},
Fig.~\ref{fig:morphology}) is not explained by differences in
cellular water activity, which are negligible between viable
hyphae of the two genera.  Instead, the genus-level variation in
$\delta$ maps onto a difference in the characteristic porosity of
each colony's aerial architecture.

\textit{Aspergillus} builds capitate conidiophores: 30--50~\textmu m
globose vesicles bearing radially packed chains of 2--3~\textmu m
conidia.  Adjacent heads abut to form a dense carpet whose porosity
spans three decades of length scale---sub-micron gaps between conidia
in a chain, few-micron gaps between adjacent chains, and
tens-of-microns gaps between neighboring heads.
\textit{Mucor} builds a sparse forest of long sporangiophores
(9--35~\textmu m diameter) each carrying a single terminal
sporangium 50--300~\textmu m in diameter; its porosity is
effectively single-scale and operates at tens to hundreds of
microns.

This geometric difference has a thermodynamic consequence through
the Kelvin equation.  For a concave water meniscus of radius $r_p$
in a wetting pore at $T = 280$~K,
$\mathrm{RH}_\mathrm{eq}(r_p) =
  \exp\!\bigl(-\ell_K / r_p\bigr)$
with Kelvin length
$\ell_K = 2\gamma V_m / (RT) \approx 1.17$~nm.
A 200~nm inter-conidial capillary depresses equilibrium relative
humidity by ${\sim}0.6\%$ below bulk; a 10~\textmu m
inter-hyphal gap by ${\sim}0.01\%$.  Because the depletion flux
near saturation scales with $(1 - \mathrm{RH})$, this sub-percent
difference is amplified into a measurable difference in sink
strength.

The multi-scale geometry analysis in Fig.~\ref{fig:morphology}
substantiates this picture.  The shallower spectral
slope~$\alpha$ of \textit{Aspergillus}
(Fig.~\ref{fig:morphology}C,\,D) reflects greater power at high
spatial frequencies and therefore greater effective surface area
per unit projected footprint.  The elevated local density CV
(Fig.~\ref{fig:morphology}E) is the statistical signature of the
dense-tuft/void texture that produces sub-micron gaps.  The
\emph{inversion} of the lacunarity comparison at the macro scale
(Fig.~\ref{fig:morphology}F)---where \textit{Mucor} is more
gap-dominated---is not contradictory but the expected signature of
\textbf{hierarchical porosity}: \textit{Aspergillus}'s closely
spaced heads fill the projected area uniformly (low lacunarity)
while remaining internally heterogeneous (high CV).  This is the
geometry required to present a large Kelvin-relevant internal
surface area without sacrificing vapor access.

We therefore interpret the fitted effective water activity
$a_{w,\mathrm{eff}}$ of the fungal sinks as a composite quantity
$a_{w,\mathrm{eff}} \sim a_{w,\mathrm{cell}} \cdot
  f_\mathrm{Kelvin}(r_p) \cdot f_\mathrm{MT}(\text{geometry})$,
in which the genus-level differences reside entirely in the two
geometric factors.  The rougher, hierarchically porous
\textit{Aspergillus} morphology simultaneously depresses
$f_\mathrm{Kelvin}$ (sub-micron capillaries) and decreases
effective mass-transfer resistance $f_\mathrm{MT}$ (higher
specific surface area), consistent with its larger $\delta$.

One caveat merits explicit statement.  The outermost layer of
\textit{A.\ fumigatus} conidia carries a RodA hydrophobin rodlet
array, with DHN-melanin in the underlying wall.  An entirely
hydrophobic conidial surface would disfavor capillary condensation.
However, the hyphal mat, conidiophore stipes, and vesicles are
hydrophilic and constitute the majority of the internal
interfacial area, and composite surfaces with hydrophobic
asperities over hydrophilic porous networks are an established
motif for enhanced dropwise condensation.  Disentangling the
geometric and chemical contributions would require comparison with
a hydrophobin-deficient mutant ($\Delta$\textit{rodA}) and is a
natural target for follow-up work.


% ============================================================
% BLOCK 8 — STATISTICAL ANALYSIS: updated sample sizes
% ============================================================

% Replace the old FFT transect description with:
Colony-surface texture was compared using automated tile-based 2D
FFT spectral slope analysis on two independent datasets
(colony surface: 14~\textit{Aspergillus}, 7~\textit{Mucor}; 3D
macro: 13~\textit{Aspergillus}, 10~\textit{Mucor} ROIs).  Light
microscopy metrics were tested across 6~\textit{Aspergillus} and
3~\textit{Mucor} fields with Benjamini--Hochberg correction for
multiple comparisons.  Surface fragmentation and lacunarity were
computed on 13~\textit{Aspergillus} and 11~\textit{Mucor} 3D
colony ROIs.  Welch's $t$-test was used throughout; Mann--Whitney
$U$~tests are reported as non-parametric confirmation.  Effect
sizes are reported as Cohen's $d$ with pooled standard deviation.
\end{antml:parameter>
</invoke>